\begin{pastexam}[label={px:ncu-114-dsa-17}]{中央資工所 114 學年度 第 17 題}
（複選；每題 5 分，全部答對才給分，答錯倒扣 1 分）

Choose correct statement(s):
\begin{enumerate}[label=(\alph*)]
\item If a problem X can polynomially reduce to an NP-hard problem, then X is NP-hard.
\item If we can prove that an existing NP-complete problem can polynomially reduce to a problem X, then
X is NP-hard.
\item An NP-hard problem can polynomially reduce to any NP-hard problem.
\item All NP problems can polynomially reduce to any NP-complete problem.
\item All NP problems can polynomially reduce to any NP-hard problem.
\end{enumerate}
\end{pastexam}

\begin{pastsoln}
答案：(b)(d)(e)。

\begin{itemize}
\item (a) 錯。 $X \le_{\mathrm{p}} Y$ 只說明 $X$ 不比 $Y$ 難，推不出 $X$ 很難。 反例：$X$ 是「給一個整數，
它是不是偶數」，它在 P 裡；歸約可以先在多項式時間內判斷答案，是就輸出一個固定能滿足的公式，否就輸出
一個固定不能滿足的公式，所以 $X \le_{\mathrm{p}} \textsf{SAT}$，而 \textsf{SAT} 是 NP 難的。 如果 $X$ 是
NP 難的，由資結書引理~28.2，NP 裡每個問題都在 P 裡，$\mathrm{P} = \mathrm{NP}$；所以除非
$\mathrm{P} = \mathrm{NP}$，$X$ 不是 NP 難的。
\item (b) 對。 這是資結書引理~28.6 證明的前半：設 $A$ 是 NP 完備的、$A \le_{\mathrm{p}} X$。 任取
$C \in \mathrm{NP}$，$A$ 是 NP 難的，所以 $C \le_{\mathrm{p}} A$；由遞移性（資結書引理~28.3），
$C \le_{\mathrm{p}} X$。 所以 $X$ 是 NP 難的。 （要說 $X$ 是 NP 完備，還要另外證 $X \in \mathrm{NP}$。）
\item (c) 錯。 資結書第~28.4~節說明了停機問題是 NP 難的，\textsf{SAT} 也是 NP 難的。 如果停機問題能
歸約到 \textsf{SAT}，先做歸約，再把 $2^n$ 組指派全部試一遍，就得到一個一定會停、而且答對停機問題的
演算法；但停機問題沒有演算法能解，矛盾。 所以停機問題不能歸約到 \textsf{SAT}。
\item (d) 對。 NP 完備的問題是 NP 難的，照定義（資結書定義~28.4），NP 裡每個問題都能歸約到它。
\item (e) 對。 這就是 NP 難的定義。
\end{itemize}
\end{pastsoln}
